\documentclass[11pt]{article}

\usepackage[a4paper,margin=25mm]{geometry}
\usepackage[T1]{fontenc}
\usepackage[utf8]{inputenc}
\usepackage{lmodern}
\usepackage{microtype}
\usepackage{amsmath,amssymb,mathtools}
\usepackage{amsthm}
\usepackage{array}
\usepackage{booktabs}
\usepackage{enumitem}
\usepackage{tabularx}
\usepackage{xcolor}
\usepackage{hyperref}

\definecolor{artifactblue}{HTML}{1F4E79}
\definecolor{artifactgreen}{HTML}{1E6B45}
\definecolor{artifactred}{HTML}{8B2F2F}
\definecolor{artifactgray}{HTML}{F2F4F5}

\hypersetup{
  colorlinks=true,
  linkcolor=artifactblue,
  urlcolor=artifactblue,
  pdftitle={A Mathematical Map of the HAETAE EasyCrypt Refinement Layer},
  pdfsubject={Specifications, refinements, and their role in HAETAE security}
}

\setlist[itemize]{leftmargin=1.5em,itemsep=0.25em,topsep=0.35em}
\setlist[enumerate]{leftmargin=1.7em,itemsep=0.25em,topsep=0.35em}
\renewcommand{\arraystretch}{1.18}
\setlength{\emergencystretch}{3em}

\theoremstyle{definition}
\newtheorem{definition}{Definition}
\theoremstyle{plain}
\newtheorem{theorem}{Theorem}
\theoremstyle{remark}
\newtheorem*{remark}{Remark}

\newcommand{\code}[1]{\nolinkurl{#1}}
\newcommand{\Z}{\mathbb Z}
\newcommand{\W}[1]{\mathbb W_{#1}}
\newcommand{\uint}[1]{\langle #1\rangle}
\newcommand{\Frame}{\mathsf{Frame}}
\newcommand{\Prefix}{\mathsf{Prefix}}
\newcommand{\Centered}{\mathsf{Centered}}

\newenvironment{resultbox}
  {\begin{center}\begin{minipage}{0.94\linewidth}\color{artifactgreen}\hrule\vspace{0.6em}\color{black}}
  {\vspace{0.6em}\color{artifactgreen}\hrule\end{minipage}\end{center}}

\newenvironment{boundarybox}
  {\begin{center}\begin{minipage}{0.94\linewidth}\color{artifactred}\hrule\vspace{0.6em}\color{black}}
  {\vspace{0.6em}\color{artifactred}\hrule\end{minipage}\end{center}}

\title{From Extracted Loops to HAETAE Security\\
\large A mathematical map of \code{spec/*.ec} and \code{refinement/*.ec}}
\author{HAETAE reference EasyCrypt artifact}
\date{Artifact state checked 2026-07-23}

\begin{document}
\maketitle

\begin{abstract}
The fifteen authored EasyCrypt files comprise six matched
specification/refinement pairs and a mode-2-parent specification with two
refinements.  They verify the generated Keccak-f transition, finite SHAKE
streams, seed-only 128-byte expansion and slices, uniform and eta
decoder/range/frame contracts, three caller schedules, and transfer of the
selected procedures into the actual-parent target.  A proof-only \(k=2,m=3\)
prefix composes the first-attempt sampler semantics and terminates with
probability one under thirteen explicit deterministic progress certificates.
This note maps those theorems into one picture.  They do not prove universal
certificate existence, semantics of \code{_keypair_full_m23}, or that the
extracted implementation realizes the probability distributions assumed by
the HAETAE EUF-CMA reduction.
\end{abstract}

\begin{resultbox}
\textbf{The result in one sentence.}
The current refinement layer establishes several necessary links from
extracted Jasmin code to well-formed sampler behavior, but the links from
SHAKE streams to ideal distributions, from the checked proof-only sampler
prefix to full key generation, and from the implementation to the abstract
security games remain separate proof obligations.
\end{resultbox}

\section{Six Pairs and One Parent Layer in Five Mathematical Notes}

The source files divide naturally by a shared prefix:

\begin{center}
\begin{tabularx}{\textwidth}{>{\raggedright\arraybackslash}p{0.15\textwidth}
  >{\raggedright\arraybackslash}p{0.40\textwidth}X}
\toprule
Layer & Authored source files & Mathematical note \\
\midrule
SHAKE256 framing &
\code{KeygenShakeStreamSpec.ec}\newline
\code{TargetKeygenShakeStream.ec} &
\emph{SHAKE256 state construction} \\
Keccak-f transition &
\code{KeygenKeccak1600Spec.ec}\newline
\code{TargetKeygenKeccak1600.ec} &
\emph{SHAKE streams and seed expansion} \\
Seed-only XOF &
\code{KeygenSeedXofSpec.ec}\newline
\code{TargetKeygenSeedXof.ec} &
\emph{SHAKE streams and seed expansion} \\
Uniform leaves &
\code{KeygenUniformXofLeafSpec.ec}\newline
\code{TargetKeygenUniformXofLeaf.ec} &
\emph{Uniform XOF leaves} \\
Eta consumer &
\code{KeygenEtaSamplerSpec.ec}\newline
\code{TargetKeygenEtaSampler.ec} &
\emph{Eta rejection sampler} \\
Caller schedules &
\code{KeygenSamplerCallersSpec.ec}\newline
\code{TargetKeygenSamplerCallers.ec} &
\emph{Sampler caller schedules} \\
Mode-2 parent prefix &
\code{KeygenMode2ParentSpec.ec}\newline
\code{TargetKeygenMode2Parent.ec}\newline
{\scriptsize\code{TargetKeygenMode2ParentComposition.ec}} &
\emph{Mode-2 actual-parent sampler prefix} \\
\bottomrule
\end{tabularx}
\end{center}

Each pair gives a vocabulary in its \code{Spec} file and proves facts about
procedures in the generated target in its \code{Target} refinement file.  A
pair is therefore best read as one mathematical unit.

\section{How to Read the EasyCrypt Judgments}

Let
\[
  \W n=\Z/2^n\Z,
\]
and write $\uint w\in\{0,\ldots,2^n-1\}$ for the unsigned representative of
$w\in\W n$.  The generated program uses words in $\W8$, $\W{32}$, and
$\W{64}$, while the specifications use ordinary integers whenever loop bounds
and array positions can be shown not to wrap.

An EasyCrypt type named \code{BArray8192} is an array of 8192 \emph{bytes}.
When accessed through 32-bit words it has $8192/4=2048$ word slots.  This unit
conversion explains why a procedure whose name ends in \code{2048} operates
on \code{BArray8192}; similarly, the \code{8192} specialization uses a
32768-byte array.

Two judgment forms occur in these files:

\begin{center}
\begin{tabularx}{\textwidth}{>{\raggedright\arraybackslash}p{0.25\textwidth}X}
\toprule
EasyCrypt form & Mathematical reading \\
\midrule
\code{hoare[P : pre ==> post]} &
Partial correctness: if the deterministic procedure $P$ is started in a state
satisfying the precondition and returns, its result satisfies the
postcondition.  A separate losslessness theorem would be needed to conclude
termination on every input. \\
\code{equiv[P ~ Q : pre ==> post]} &
Relational equivalence: executions of $P$ and $Q$ are compared from related
inputs.  In the caller file the exported postcondition mentions only equality
of returned arrays, not equality of an exposed call trace. \\
\bottomrule
\end{tabularx}
\end{center}

No procedure in this fifteen-file surface contains an EasyCrypt probabilistic
sampling command.  Distributional conclusions therefore require additional
lemmas connecting deterministic XOF output to an abstract probabilistic model.

\section{The Verified Layers}

\subsection{Exact Keccak, SHAKE, and seed expansion}

Starting from a zero 200-byte Keccak state, the checked initializer absorbs
64 selected seed bytes and the two least-significant nonce bytes, then inserts
the SHAKE suffix and the final rate bit.  Expanding the proved recurrence gives
the rate-block byte string
\[
  S[o:o+64]\;\Vert\;\operatorname{LE}_2(\nu)\;\Vert\;
  \mathtt{1f}\;\Vert\;0^{68}\;\Vert\;\mathtt{80},
\]
under the premise $\uint o+64\le128$.  Companion theorems prove the generated
24-round Keccak-f transition, exact rate-block serialization, and bounded
consecutive SHAKE128/SHAKE256 squeeze blocks.  The seed-only Hoare theorem
proves on return that \code{_kp_expand_seedbuf} equals the first 128 SHAKE256
bytes of its 32-byte seed; named corollaries identify output ranges $0$--$31$,
$32$--$95$, and $96$--$127$.

\subsection{Uniform leaf range and frame}

Put $q=64513$ and $N=256$.  For either destination capacity
$m\in\{2048,8192\}$ word slots, let an extracted uniform leaf return $A'$ from
input array $A$ at word base $b$.  Under $b+N\le m$, the two checked
postconditions are
\[
  0\le \operatorname{uint}_{32}(A'[b+i])<q
  \qquad(0\le i<N)
\]
and
\[
  A'[j]=A[j]
  \quad\text{for every byte }j\notin[4b,4(b+N)).
\]
They are conditional on return.  Internally, a finite-buffer consumer forms a
16-bit little-endian candidate and accepts it exactly when it is below $q$.
Stronger stream theorems connect the four-block prelude and every later block
to one exact finite SHAKE128 byte sequence on return.  Both bounded generated
consumers terminate unconditionally.  That fact does not prove that either
whole uniform leaf eventually obtains 256 accepted candidates.  Separately,
\code{uniform2048_leaf_progress_ll} and
\code{uniform8192_leaf_progress_ll} prove probability-one termination of the
two actual leaves under exact seed, slice, base, and capacity premises plus a
concrete \code{uniform_progress_prefix} certificate.  Its finite endpoint has
at least 256 accepted candidates, and every preceding still-incomplete
full-block prefix from block four makes strict progress.  The certificate is
not proved for every input.  The parent refinement lifts a six-certificate
bundle through the fixed \(2\times3\) matrix caller and a two-certificate
bundle through the fixed vector caller.

\subsection{Eta consumer support, counter, and frame}

For a 32-bit word $x$, define the centered residue
\[
  \chi(x)=
  \begin{cases}
    -1,&\operatorname{uint}_{32}(x)\bmod3=2,\\
    \operatorname{uint}_{32}(x)\bmod3,&\text{otherwise}.
  \end{cases}
\]
The generated arithmetic helpers are proved to compute the word encoding of
$\chi(x)$ on their stated input ranges.  The key identity is
\[
  \left\lfloor\frac{171x}{512}\right\rfloor
  =\left\lfloor\frac{x}{3}\right\rfloor,
  \qquad 0\le x\le242.
\]

The finite eta consumer rejects a byte $x\ge243=3^5$.  From an accepted byte
it attempts to write the five centered base-3 digits, stopping early if the
256-word destination is full.  If the initial counter is at most 256 and any
existing prefix is already centered, then on return:
\begin{itemize}
\item the counter has not decreased and is at most 256;
\item every word in the returned established prefix has signed value in
  $\{-1,0,1\}$; and
\item no word outside the reserved 256-word block has changed.
\end{itemize}
The current file also composes the exact capped decoder through successive
136-byte SHAKE256 overwrite blocks.  On return, the whole eta leaf has exactly
256 centered coefficients, and the actual eta caller carries a separate exact
finite witness for every selected vector slot.  The bounded eta consumer is
unconditionally lossless.  In addition, \code{eta2048_leaf_progress_ll} proves
probability-one termination of the actual eta leaf under exact seed, slice,
base, and capacity premises plus a concrete
\code{eta_progress_prefix} certificate.  The certificate supplies a finite
full endpoint and strict decoded-count progress at each preceding incomplete
block prefix.  Its existence is not proved for every input.  The parent
refinement lifts five such certificates through the two fixed eta callers of
the first attempt; it does not prove totality of the actual retry loop.

\subsection{Caller schedules}

The caller refinement treats the three sampler leaves as fixed opaque partial
maps.  It proves that the extracted matrix, uniform-vector, and eta-vector
loops return the same arrays as transparent integer-indexed schedules.  The
essential bases and nonces are
\[
\begin{array}{c|c|c}
\text{caller} & \text{word base} & \text{nonce}\\
\hline
\text{matrix }(i,j) & [256(ic+j)]_{64} & [256i+j]_{64}\\
\text{uniform vector }i & [256i]_{64} & [256k+m+i]_{64}\\
\text{eta vector }i & [256i]_{64} & \nu+[i]_{64}.
\end{array}
\]
This relational result is an iterator-correctness theorem.  Separate Hoare
theorems supply finite leaf semantics, ranges, and aggregate caller frames.
Seven cross-module equivalences transport the seed, three leaf, and three
caller results to procedures in the actual-parent target.  An authored
proof-only prefix then composes seed expansion, the \(2\times3\) matrix caller,
the two-slot uniform-vector caller, and eta callers for nonces \(0\)--\(4\).
Its semantic theorem is conditional on return; a separate pHoare theorem
proves probability-one termination under the bundled \(6+2+5\) progress
certificates.  Neither theorem is about \code{_keypair_full_m23}.

\section{What Can and Cannot Be Composed Today}

The local layers are deliberately compatible: the caller theorem invokes the
same fixed uniform and eta leaves studied by the local Hoare proofs.  The
parent layer now crosses one compositional cut, through the first sampler
attempt of a proof-only \(k=2,m=3\) prefix.  No end-to-end key-generation
theorem is exported, however.

\begin{center}
\setlength{\fboxsep}{7pt}
\begin{tabular}{c}
\fcolorbox{artifactgreen}{artifactgray}{\parbox{0.78\textwidth}{\centering
Generated EasyCrypt targets from the sampler-caller and actual-parent
Jasmin closures}}\\[0.6em]
$\Downarrow$ \quad extraction gates, six proof pairs, and parent triplet\\[0.6em]
\fcolorbox{artifactgreen}{artifactgray}{\parbox{0.78\textwidth}{\centering
Exact Keccak/SHAKE bytes and seed slices; uniform/eta finite decoders;
range/frame; bounded-consumer totality; certificate-conditioned uniform/eta
leaf totality; caller schedules}}\\[0.6em]
$\Downarrow$ \quad proof-only actual-parent sampler prefix checked\\[0.6em]
\fcolorbox{artifactgreen}{artifactgray}{\parbox{0.78\textwidth}{\centering
Parent-qualified seed/caller semantics and first-attempt sampler-prefix
composition for \(k=2,m=3\), with certificate-conditioned totality}}\\[0.6em]
$\Downarrow$ \quad \textcolor{artifactred}{missing distributions,
\code{_keypair_full_m23}, retry/downstream semantics, and packing}\\[0.6em]
\fcolorbox{artifactred}{artifactgray}{\parbox{0.78\textwidth}{\centering
Abstract HAETAE key generation with the distributions assumed by the scheme}}\\[0.6em]
$\Downarrow$ \quad \textcolor{artifactred}{missing keygen/sign/verify
implementation refinement}\\[0.6em]
\fcolorbox{artifactred}{artifactgray}{\parbox{0.78\textwidth}{\centering
Implementation-level EUF-CMA theorem}}
\end{tabular}
\end{center}

The green boxes are supported by the extraction and proof gates.
The red boxes name composition goals, not consequences of the present
theorems.

\section{Connection to the Provable-Security Development}

The repository's separate provable-security development contains an abstract
HAETAE scheme and EUF-CMA bounds.  For example,
\code{HAETAE_Security.ec} bounds the forgery probability under explicit
game-hop and hardness premises.  That theorem concerns the abstract module
\code{HAETAE}; it is not automatically a theorem about the Jasmin program.

An implementation-level security statement needs a bridge of the following
logical form.  If an implementation $I$ and abstract scheme $H$ have the same
observable behavior in the security experiment, then
\[
  \Pr[\operatorname{Forge}(I)]
  =\Pr[\operatorname{Forge}(H)]
  \le \operatorname{HAETAE\_EUF\_bound}.
\]
The files documented here contribute local premises for the first equality;
they do not prove it.  The repository traceability table accordingly records
the sampler-XOF path as partial and the paper-faithful EUF-CMA surface as
partial.  Full mode-2 Jasmin security composition remains unchecked.

\begin{boundarybox}
\textbf{Security interpretation.}
Range and frame postconditions are safety properties.  They prevent malformed
coefficients and out-of-region functional-array writes on terminating runs.
They do not imply that coefficients are uniform, independent, unpredictable,
or distributed as required by a security reduction.
\end{boundarybox}

\section{Remaining Proof Obligations}

A complete bridge from these local refinements to HAETAE provable security
still requires, at minimum:

\begin{enumerate}
\item semantic composition with the actual \code{_keypair_full_m23} past the
  current proof cut, including matrix/NTT work, finalize, singular rejection,
  retry control, seed selection, and packing;
\item probabilistic theorems for the uniform rejection sampler and centered
  base-3 sampler, plus universal progress-certificate existence or another
  unconditional almost-sure-termination argument;
\item certificates and semantic contracts for later retry attempts, together
  with an acceptance or retry-termination argument;
\item refinement of the complete key-generation result to the abstract HAETAE
  key-generation distribution, including remaining adjustments documented by
  the repository;
\item corresponding implementation refinements for signing and verification;
  and
\item a final game-level equivalence transferring the abstract EUF-CMA bound
  to the target implementation, together with separate memory-safety and
  leakage claims if those are desired.
\end{enumerate}

This list separates functional correctness from cryptographic assumptions:
proving a correct deterministic decoder is different from proving that its
input is an ideal or pseudorandom XOF stream.

\section{Suggested Reading Order}

For a mathematical reader, the most efficient order is:

\begin{enumerate}
\item \code{TargetKeygenSamplerCallers.pdf}: the simplest proof, consisting of
  closed-form schedules and ordinary induction;
\item \code{TargetKeygenUniformXofLeaf.pdf}: frame and range invariants for a
  rejection sampler;
\item \code{TargetKeygenEtaSampler.pdf}: centered base-3 arithmetic and the
  five-write consumer invariant; and
\item \code{TargetKeygenShakeStream.pdf}: exact byte/lane framing at the XOF
  boundary; and
\item \code{TargetKeygenMode2Parent.pdf}: transfer into the actual-parent
  target and the checked first-attempt sampler prefix.
\end{enumerate}

Return to the present map after reading the five notes to see which arrows
toward the security theorem are checked and which are still open.

\section{Machine-Checked Evidence}

From the \code{haetae-ref-easycrypt/} directory, run
\begin{quote}
\begin{verbatim}
./scripts/verify-keygen-sampler-callers-proof.sh
\end{verbatim}
\end{quote}
The gate checks pinned source hashes and extraction drift across the unified
25-file/31-procedure target, then compiles one generated target followed by six
specifications and six refinements with \code{-no-eco}.  The retained
2026-07-22 full run reports 13 compiled entries and passing scans of the
authored files for proof holes and axiom declarations.  The actual mode-2 parent has a separate
32-file/56-procedure extraction gate.  The dedicated semantic gate is
\begin{quote}
\begin{verbatim}
./scripts/verify-keygen-mode2-parent-proof.sh
\end{verbatim}
\end{quote}
It checks both extraction closures and their 24 byte-identical shared
generated theories, then compiles a 17-entry \code{-no-eco} manifest containing
both targets, seven specifications, and eight refinements.  The retained
2026-07-23 run reports 17 of 17 entries compiled and passing source, drift,
shared-file identity, proof-hole, and axiom checks.

The source-to-document status map is maintained in
\code{manifests/traceability.csv}.  The abstract security development and its
own verification instructions live under
\code{haetae-security/provable-security/}.

\section*{Conclusion}

The current \code{spec/*.ec} files define the mathematical predicates needed
to reason about concrete arrays, counters, and framing.  The corresponding
\code{refinement/*.ec} files prove that selected generated procedures satisfy
those predicates.  Together they form a useful deterministic foundation for
HAETAE implementation security.  Exact finite XOF bytes are no longer the
missing bridge; universal progress certificates, generic caller totality,
sampler
distributions, full \code{_keypair_full_m23} and retry/downstream semantics,
source-pointer properties, full-algorithm refinement, and game composition
remain open.  A proof-only actual-parent sampler prefix now composes the
previous local results through the first five eta nonces.

\end{document}
